\section{实验设计与结果整理}
\label{sec:experiments}
本节围绕四个问题组织评价：语言引导能否改善最终三维分区；这种变化来自哪些提示与后处理操作；收益是否伴随额外的时间和失败风险；保存的名称是否真正对应目标部件。\textbf{当前数据状态：}所提供的三份工作簿均为待填模板，指标单元格全部为 TBD。因此，本节给出已明确的评价口径、比较方案和分析结构，不报告尚未取得的定量优越性。前期小样本目视观察只用于排查流程，不替代最终测试。

\subsection{数据来源与配对管理}
研究样本来自 Objaverse 三维资产，候选清单共 684 个对象，包括组装结构（Assemblies）220 个、实体体积结构（Solid-volume）232 个和细薄／细粒度结构（Thin/fine-grained）232 个。该划分用于分析不同几何结构上的行为，并非互斥语义类别的通用定义。对象筛选准则、授权信息、重复对象检测及人工参考制作说明应与清单一同整理；最终开发集和测试集的具体 ID 与数量为 \TBD{数据划分}。已经用于提示调整或故障分析的对象应归入开发集，不再作为未见测试样本支持最终结论。

每个对象保留原始下载资产、相机参数、RGBA 图像、生成资产、VXZ、二维标签图和推理配置之间的对应关系。图像一致性包括尺寸、alpha、视角、方向以及内容指纹，不能只比较前景轮廓。D 参考来自人工直接制作和渲染，不以 C 的预测或语言提议作为唯一标注依据。对于无法判断的遮挡边界，可设置忽略区并单独统计，而不是强制赋予一个确定标签。

人工参考首先定义当前视角的部件划分；若生成几何与原始资产轮廓发生变化，源视角参考未必能直接评价新表面。计算投影指标前需要核验坐标变换、相机和轮廓对齐，并说明参考是否对应被评价的生成几何。对齐不成立的样本不得仅靠缩放、重涂或最近邻颜色匹配强行取得分数，应列为不可评价项并报告数量。

质量比较使用预先确定的配对测试集合，同时报告每种方法的尝试数、成功数和可评价数。成功交集上的均值只反映条件于成功的质量，不能掩盖失败差异；失败率另以全部测试尝试为分母。结构分组先计算对象级指标，再报告各组均值及三组等权宏平均。全体对象加权均值可作为补充，但不得与等权宏平均混用。

\subsection{对照方法与条件公平性}
实验分为两层。第一层在冻结的 SegViGen 下比较 A/B/C/D 条件，隔离二维引导的作用。A 为原始 RGBA，B 为自动 SAM2，C 为本文 V5 语言规划、反馈和融合，D 为人工二维参考条件。D 既用于构造二维参考，也可作为三维推理的人工引导对照；但不将 D 对自身的二维评分作为普通方法结果，更不将 D 条件的三维输出直接称为三维真值。该层对比还需补充 A/B/C/D 的最终配对明细，不能由外部方法表自动推导。

第二层比较不同三维分割机制，方法列于表\ref{tab:main}。Find3D\citep{ma2024find3d}使用预先固定来源的部件名称；PartField\citep{liu2025partfield}采用其无部件提示设置；P3-SAM\citep{ma2025p3sam}和 PartSAM\citep{zhu2025partsam}采用可复现的自动分解配置；SegviGen\citep{li2026segvigen}采用无部件提示的对照设置。所有方法处理同一生成几何，记录输入采样、粒度和聚类等设置，不能按测试参考为某一方法逐对象选择最佳部件数。

GeoSAM2\citep{deng2025geosam2}的原生接口接受二维空间提示并使用多视图几何信息。工作簿计划为其提供与本文同源的二维引导；执行时需要明确如何把共享引导转换为该接口支持的提示，以及是否包含额外视图。该设置评价的是“共享提示＋不同三维处理”的系统，而非未经适配的默认模型。输入条件不同的方法并列展示时，需保留条件列并说明额外语义信息来源。D 参考及测试标签不能用于生成仅供某个基线享有的理想提示。

\begin{table*}[t]
\caption{生成资产上的最终三维部件分割比较。保留所提供工作簿的全部方法与指标；所有测量值目前均为 TBD。投影指标评价最终三维输出在固定源视角下与人工二维参考的一致性，不代表完整三维准确率。部件合理性和边界质量为拟开展的匿名多视图人工评分。}
\label{tab:main}
\centering\small
\setlength{\tabcolsep}{5pt}
\begin{tabularx}{\textwidth}{@{}l X c c c c@{}}
\toprule
方法 & 部件条件 & \shortstack{投影 mIoU\,(\%) $\uparrow$} & \shortstack{投影 BF1\,(\%) $\uparrow$} & \shortstack{部件合理性\,(1--5) $\uparrow$} & \shortstack{边界质量\,(1--5) $\uparrow$}\\
\midrule
Find3D & 部件名称 & TBD & TBD & TBD & TBD\\
PartField & 无部件提示 & TBD & TBD & TBD & TBD\\
P3-SAM & 自动提示 & TBD & TBD & TBD & TBD\\
PartSAM & 自动提示 & TBD & TBD & TBD & TBD\\
GeoSAM2 & 共享二维引导 & TBD & TBD & TBD & TBD\\
SegviGen & 无部件提示 & TBD & TBD & TBD & TBD\\
本文方法 & 生成二维引导 & TBD & TBD & TBD & TBD\\
\bottomrule
\end{tabularx}
\vspace{2pt}
\begin{minipage}{\textwidth}\footnotesize
注：GeoSAM2 的共享引导是拟议的受控设置，须通过其支持的提示接口实现并记录适配方法；不可默认其原生接受任意彩色分割图。Find3D 名称来源、各方法粒度配置、测试子集、有效输出数及配对区间均待冻结。本文不使用二维引导分数代替最终三维输出分数。
\end{minipage}
\end{table*}


\subsection{二维与投影评价指标}
\paragraph{二维匹配 mIoU.}
预测部件集合记为 $\mathcal{P}=\{P_i\}_{i=1}^{m}$，参考集合为 $\mathcal{D}=\{D_j\}_{j=1}^{n}$。剔除预定义忽略像素，并将背景 ID 排除在待匹配部件集合之外，计算
\begin{equation}
 J_{ij}=\frac{|P_i\cap D_j|}{|P_i\cup D_j|}.
\end{equation}
使用一对一匹配最大化 IoU 总和，未匹配参考按零计，定义
\begin{equation}
 \mathrm{mIoU}_{\rm match}
 =\frac{1}{n}\max_{\pi}\sum_{j\in\operatorname{dom}(\pi)}J_{\pi(j),j}.
 \label{eq:miou}
\end{equation}
此处 $\pi$ 是参考到预测的部分单射。预测泄漏到参考背景上的像素仍进入并集并受到惩罚。该实现以参考部件数为分母，而非仅对成功匹配的部件求均值，因此漏件不能被排除。多余预测不直接进入该分母，故同时报告部件 F1 约束过分割。相同整数标签的分离区域仍视为同一预测部件，不能在评价时按连通域重新拆分。

\paragraph{边界 F1.}
对 IoU 最优匹配中重叠非零的每一对区域提取内边界，在距离容差 $\delta$ 内计算边界精确率与召回率，再求调和平均。按参考部件数平均，未匹配参考及 IoU 为零的匹配对均计零。现有评价代码的默认值为 $\delta=2$ 像素，并排除忽略区邻近边界；最终实验需固定分辨率与容差。不同缩放分辨率下的同一个像素阈值代表不同几何容差，不能直接混合报告。

\paragraph{部件 F1@0.5.}
在 $J_{ij}\geq0.5$ 的候选边上寻找最大基数的一对一匹配，以匹配数为 $\mathrm{TP}$，定义
\begin{equation}
 \mathrm{PartF1}_{0.5}=\frac{2\mathrm{TP}}{m+n}.
 \label{eq:partf1}
\end{equation}
这一匹配目标是最大化达到阈值的部件数，与式\eqref{eq:miou}的最大 IoU 和匹配不同。该指标同时惩罚漏检和额外预测；单纯比较输出部件数量，既不能证明召回提高，也不能证明过分割减少。

\paragraph{最终三维输出的投影指标.}
将最终分割 GLB 在固定源相机下以不受光照影响的标签方式渲染为 $\widehat L^{\rm proj}$，再与人工二维参考计算上述匹配和边界指标。若只能读取颜色编码输出，应固定调色板解码、颜色空间、背景和不确定像素处理，不将带光照的 RGB 差异当作语义差异。投影过程必须使用最终三维输出，不得复用其输入的 C 引导图。该评价只覆盖可见表面，不能说明背面正确性、完整部件几何或跨视角一致性。

\paragraph{人工三维评价与名称核验.}
计划对随机顺序、匿名展示的多视图或转台结果分别评估部件合理性和边界质量，均使用 1--5 分。前者关注划分是否对应可解释结构、是否存在明显漏件和过细碎片，后者关注边界是否沿部件交界、是否粘连或跨越不同结构。评分者人数、训练示例、分值锚点、分歧处理和一致性统计为 \TBD{人工评价协议}。名称核验则以保留部件的掩码及英文词组为对象，分别记录定位匹配、语义正确和重复命名；字段非空不等于标签正确。三维点击命名只有在表面到 ID 的映射完成后才另行评价。

\subsection{主实验的分析结构}
表\ref{tab:main}预留最终三维输出的比较结果。填入测量值后，首先检查投影与人工多视图评价是否给出一致判断：投影改善而背面表现下降，可能意味着源视角引导未能稳定传播；合理性提高但边界评分不变，可能反映语义粒度改善而非几何边界精修。反之，较高的边界分数也不必然意味着名称语义正确。

对 B 与 C 的比较，应同时检查匹配 mIoU、部件 F1 和错误类型。B 输出更多区域时，增加的区域可能是真实小件，也可能是重复纹理或碎片；C 输出更少区域时，既可能是合理合并，也可能是漏件。本文不根据区域数预先选择解释，而用相同参考、局部放大和人工核验区分这些情况。当前没有足够的最终测量值判定哪种现象占主导，主结论将在配对结果完成后填写。

\subsection{组件消融与提示预算}
表\ref{tab:ablation}保留工作簿定义的七种配置。自动 SAM2 用于观察无语言规划的结果；框＋点基线用于比较较简化的空间提示策略；其余变体从完整方法中分别限制正点数量、删除负点、关闭反馈或替换点感知融合。每次改变需固定图像、模型、种子、三维资产及其余处理，保存实际送入 SAM2 的点坐标与正负标签，而不只记录配置名称。

对于点感知融合的消融，需要在运行前确定替代方案。当前模块同时包含点连通清理、小框优先排序和残余填充；若将它整体替换为普通覆盖，只能归因于组合模块，不能宣称单独证明某一规则有效。若要解释“小件保护”的机制，应追加只改覆盖顺序的严格对照，并统计按参考面积划分的小、中、大部件表现。面积分组阈值应由开发集预先确定，不按测试结果调整。

反馈实验还应报告 SAM2 重试触发比例、重试成功率及单对象额外耗时，并区分“反馈判定合格”“修正后保留”“反馈失败回退”。完整方法对每个初始有效候选都发起反馈查询，只有 SAM2 重试为条件触发，计时需包括所有查询。关闭反馈的样本与完整方法应使用相同初始规划，以免把 Qwen 采样变化误认为反馈收益。提示预算折线图需要独立控制点数和负点策略；不能把不同提示、不同融合、不同反馈设置的行连接起来解释为单变量趋势。

\begin{table*}[t]
\caption{空间提示、反馈和融合的消融实验。前三项评价三维推理之前的二维引导，最后一项评价最终三维输出的固定视角投影。前两行为比较基线，随后四行为相对完整方法的组件消融，最后一行为完整方法；全表尚无测量结果。}
\label{tab:ablation}
\centering\small
\setlength{\tabcolsep}{7pt}
\begin{tabularx}{\textwidth}{@{}X c c c c@{}}
\toprule
配置 & \shortstack{二维匹配 mIoU\,(\%) $\uparrow$} & \shortstack{二维 BF1\,(\%) $\uparrow$} & \shortstack{二维部件 F1@0.5\,(\%) $\uparrow$} & \shortstack{投影 mIoU\,(\%) $\uparrow$}\\
\midrule
Automatic SAM2 & TBD & TBD & TBD & TBD\\
Qwen--SAM2（框＋点） & TBD & TBD & TBD & TBD\\
本文方法：仅单个正点 & TBD & TBD & TBD & TBD\\
本文方法：移除负点 & TBD & TBD & TBD & TBD\\
本文方法：移除反馈修正 & TBD & TBD & TBD & TBD\\
本文方法：移除点感知融合 & TBD & TBD & TBD & TBD\\
本文完整方法 & TBD & TBD & TBD & TBD\\
\bottomrule
\end{tabularx}
\vspace{2pt}
\begin{minipage}{\textwidth}\footnotesize
注：“仅单个正点”保留框、负点、反馈及融合，其与独立框＋点基线的差异必须由配置清单说明。移除点感知融合时的替代规则尚需固定；若同时删除多个操作，结果只能解释为该组合模块的贡献。匹配基于区域重叠，不按调色板 RGB 值一一对应。
\end{minipage}
\end{table*}


\subsection{效率、可靠性与失败案例}
表\ref{tab:efficiency}比较端到端成本。计时包含方法特有的预处理、提示构造、分割和后处理。本文方法需计入 Qwen、SAM2、反馈和 SegViGen；如果多个阶段顺序执行，峰值显存取全过程实际观测最大值，而非简单相加各阶段峰值。共享 TRELLIS.2 资产生成与格式转换成本另列，并说明离线缓存是否计入。硬件型号、软件版本、模型加载是否计时、预热次数和重复次数均为 \TBD{运行协议}。

运行失败包括无法完成所需阶段或无法导出合法结果，其比例以所有测试对象尝试为分母。分割坍塌则指已导出结果中的退化划分，与程序崩溃分开统计；其判据需要在测试前定义，不能把所有单部件对象都认作失败。最终报告应同时列出分子、分母、重试次数和缺失原因，避免仅展示成功子集的时间或质量。

定性分析将覆盖三类结构，并保留漏件、部件粘连、碎片、重复候选、错误负点和语义误名等失败类型。图\ref{fig:qualitative}至图\ref{fig:budget-lines}暂留空白。图像必须来自同一次冻结实验，不将旧 V4 的引导与 V5 的三维输出配成一行；方法与参考的视角、缩放和配色语义应保持可比较。待对象级结果完整后，计划以对象为单位进行配对重采样，报告差值及置信区间，而非仅凭单个最佳样本宣称有效。

\begin{table*}[t]
\caption{效率与可靠性。端到端耗时包含各方法所需的预处理、提示生成、分割与后处理；共享 TRELLIS.2 资产生成成本单列。失败率以全部测试尝试为分母，坍塌率以成功导出的有效结果为分母。硬件、计时协议和坍塌判据尚待冻结，所有结果为 TBD。}
\label{tab:efficiency}
\centering\small
\begin{tabularx}{\textwidth}{@{}X c c c c@{}}
\toprule
方法 & \shortstack{总耗时\,(s/对象) $\downarrow$} & \shortstack{峰值显存\,(GB) $\downarrow$} & \shortstack{运行失败率\,(\%) $\downarrow$} & \shortstack{分割坍塌率\,(\%) $\downarrow$}\\
\midrule
Find3D & TBD & TBD & TBD & TBD\\
PartField & TBD & TBD & TBD & TBD\\
P3-SAM & TBD & TBD & TBD & TBD\\
PartSAM & TBD & TBD & TBD & TBD\\
GeoSAM2 & TBD & TBD & TBD & TBD\\
SegviGen & TBD & TBD & TBD & TBD\\
本文方法 & TBD & TBD & TBD & TBD\\
\bottomrule
\end{tabularx}
\vspace{2pt}
\begin{minipage}{\textwidth}\footnotesize
注：本文方法计入 Qwen、SAM2、反馈和 SegViGen；GeoSAM2 的共享引导成本不得漏计。若方法需要逐对象优化，其优化时间计入总耗时。缺测项不是零；历史版本的成功率不替代本版结果。
\end{minipage}
\end{table*}

\begin{figure*}[t]
  \centering
  \blankfigure{22mm}{可视化结果待补}{输入 RGBA、二维 B/C/D 引导、对应三维结果与局部放大。使用同一对象、同一相机；同时保留成功与失败样本。}
  \caption{定性结果占位。计划展示组装结构、实体体积结构与细薄结构对象的二维引导及三维分割对比；本版不放入未经配对核验的结果。}
  \label{fig:qualitative}
\end{figure*}
\begin{figure*}[t]
  \centering
  \begin{minipage}[t]{0.48\textwidth}
    \centering
    \blankfigure{18mm}{分组性能柱状图待补}{按三类对象分别汇总主指标，附对象级配对置信区间。}
    \caption{结构类别对结果的影响（占位）。}
    \label{fig:group-bars}
  \end{minipage}\hfill
  \begin{minipage}[t]{0.48\textwidth}
    \centering
    \blankfigure{18mm}{消融差值柱状图待补}{以完整方法为参照，报告同一测试子集上的指标变化。}
    \caption{提示、反馈与融合的消融分析（占位）。}
    \label{fig:ablation-bars}
  \end{minipage}
\end{figure*}
\begin{figure*}[t]
  \centering
  \blankfigure{18mm}{提示预算—质量／耗时折线图待补}{横轴使用预先定义的点提示预算；分别绘制质量与端到端时间，不预设单调改善。}
  \caption{提示预算与质量、计算成本之间的关系（占位）。本图属于追加实验，不能从表\ref{tab:ablation}的若干异质配置直接连线得到。}
  \label{fig:budget-lines}
\end{figure*}

